\documentclass[12pt, a4paper]{article}

\usepackage[utf8]{inputenc}
\usepackage[T1]{fontenc}

\usepackage{textcomp}
\usepackage{geometry}
\geometry{margin=2.5cm}
\usepackage{titlesec}
\usepackage{enumitem}
\usepackage{amsmath, amssymb}
\usepackage{hyperref}
\usepackage{xcolor}
\usepackage{mdframed}


% Colori
\definecolor{accent}{RGB}{30, 80, 140}
\definecolor{boxbg}{RGB}{240, 245, 255}
\definecolor{fontecolor}{RGB}{90, 90, 90}


% Box fonti
\newmdenv[
backgroundcolor=boxbg,
linecolor=accent,
linewidth=1pt,
leftmargin=0pt,
rightmargin=0pt,
innerleftmargin=10pt,
innerrightmargin=10pt,
innertopmargin=6pt,
innerbottommargin=6pt,
skipabove=8pt,
skipbelow=8pt
]{fontebox}

\newcommand{\fonti}[1]{%
	\begin{fontebox}
		\textcolor{fontecolor}{\footnotesize\textbf{Fonti:} #1}
	\end{fontebox}%
}

% Comando per nota concettuale
\newcommand{\nota}[1]{%
	\par\smallskip
	\noindent\textit{\small $\triangleright$\; #1}
	\par\smallskip
}

\setlist[itemize,1]{leftmargin=1.5em, itemsep=2pt, topsep=4pt}
\setlist[itemize,2]{leftmargin=1.5em, itemsep=1pt, topsep=2pt}

\hypersetup{
	colorlinks=true,
	linkcolor=accent,
	urlcolor=accent
}

\title{%
	\textbf{Strutture Hamiltoniane in Meccanica Quantistica}\\[6pt]
	\large Perché le Strutture della Meccanica Razionale Riemergono\\
	nella Meccanica Quantistica senza Essere Assunte negli Assiomi?\\[4pt]
	\normalsize Indice Concettuale --- Versione 5
}
\author{Tommaso Pedroni}
\date{\today}

\begin{document}
	
	\maketitle
	\tableofcontents
	\newpage
	
	% ============================================================
	\section*{Domanda Guida della Tesi}
	% ============================================================
	
	Le strutture della meccanica razionale --- spazio delle fasi, parentesi di Poisson,
	flusso hamiltoniano, teorema di Noether --- riemergono nella meccanica quantistica
	senza essere assunte negli assiomi. La tesi argomenta che questo non è una coincidenza
	né un'analogia formale, ma è una conseguenza strutturale del fatto che meccanica
	razionale e meccanica quantistica sono \textbf{fibre diverse dello stesso campo continuo
		di $C^*$-algebre}. La risposta si articola in tre passi di profondità crescente.
	
	\medskip
	
	\textbf{Passo 1 (Parti I--II, Capp.\ 1--5).}
	Classico e quantistico sono rappresentazioni della stessa struttura astratta: una
	$C^*$-algebra. La meccanica razionale ne è la fibra commutativa
	($A_0 = C_0(M)$, teorema di Gelfand), la meccanica quantistica ne è la fibra
	non commutativa ($A_\hbar \subset B(\mathcal{H})$, teorema di Gelfand--Naimark).
	Le strutture non riemergono perché non erano mai scomparse: parentesi di Poisson
	e commutatore sono entrambi Lie-bracket della rispettiva $C^*$-algebra.
	
	\medskip
	
	\textbf{Passo 2 (Cap.\ 6, Rieffel).}
	Non esiste un salto discreto tra i due regimi: la strict deformation quantization
	di Rieffel costruisce una famiglia continua di $C^*$-algebre $\{A_\hbar\}_{\hbar \in [0,1]}$.
	La condizione di Dirac
	$\bigl\|\frac{1}{i\hbar}[Q_\hbar(f), Q_\hbar(g)] - Q_\hbar(\{f,g\})\bigr\| \to 0$
	garantisce che le strutture classiche persistano in senso normato, non per analogia.
	
	\medskip
	
	\textbf{Passo 3 (Cap.\ 9, Ashtekar--Schilling).}
	Lo spazio degli stati quantistico $\mathbb{P}\mathcal{H}$ è esso stesso una varietà
	simplettica (kähleriana). La meccanica quantistica \emph{è} meccanica hamiltoniana su
	uno spazio degli stati diverso; non c'è nulla da spiegare, perché non c'è nulla che emerga.
	
	% ============================================================
	\newpage
	\section*{Introduzione: Struttura e Obiettivi del Lavoro}
	% ============================================================
	
	La presente tesi risponde a un interrogativo fondazionale della fisica matematica:
	per quale motivo le strutture portanti della meccanica razionale ---
	spazio delle fasi, parentesi di Poisson, flusso hamiltoniano, conservazione di Noether
	--- riemergono intatte nel formalismo della meccanica quantistica, pur non essendo mai
	postulate nei suoi assiomi?
	
	L'obiettivo è mostrare che tale corrispondenza non è un'analogia storica né un
	principio di corrispondenza assunto come ipotesi aggiuntiva, ma una rigorosa
	conseguenza strutturale. La strategia adottata è \emph{doppia}: un approccio
	algebrico \textit{bottom-up} (Parti I--II) che costruisce entrambe le teorie come
	istanze della stessa struttura di $C^*$-algebra, e un approccio geometrico
	\textit{top-down} (Parte III) che mostra come le strutture classiche siano
	letteralmente presenti nello spazio degli stati quantistico.
	
	\medskip
	
	\textbf{Il lato classico (Capp.\ 1--3).}
	La Parte~I costruisce il linguaggio geometrico e algebrico della meccanica razionale
	nel modo più invariante possibile: lo spazio delle fasi come varietà simplettica,
	le osservabili come algebra di Lie $C^\infty(M)$, il teorema di Noether come proprietà
	delle mappe momento. Questa formulazione non è un lusso tecnico: è il linguaggio
	che permetterà di riconoscere le stesse strutture nel mondo quantistico.
	
	\medskip
	
	\textbf{Il quadro algebrico (Capp.\ 4--6).}
	La Parte~II stabilisce che entrambe le teorie sono $C^*$-algebre (Cap.\ 4),
	giustifica il passaggio dagli operatori illimitati all'algebra di Weyl tramite
	il teorema di Wintner--Wielandt (Cap.\ 5), e costruisce il ponte continuo tra
	classico e quantistico tramite la strict deformation quantization di Rieffel (Cap.\ 6).
	
	\medskip
	
	\textbf{Le risposte (Capp.\ 7--10).}
	La Parte~III risponde alla domanda capitolo per capitolo: il Lie-bracket (Cap.\ 7),
	il flusso hamiltoniano (Cap.\ 8), la struttura simplettica dello spazio degli stati
	(Cap.\ 9), il teorema di Noether e il taglio di Heisenberg (Cap.\ 10).
	L'esempio del qubit attraversa tutti e dieci i capitoli come filo trasversale.
	
	% ============================================================
	\newpage
	\part*{Parte I --- Il Lato Classico}
	\addcontentsline{toc}{part}{Parte I: Il Lato Classico}
	% ============================================================
	
	\nota{La Parte~I costruisce il linguaggio geometrico della meccanica razionale in forma
		invariante. Non si usano coordinate più del necessario: tutto ciò che viene formulato
		qui sarà riconoscibile quando riapparirà nella Parte~III. L'obiettivo non è ripassare
		la meccanica classica --- è preparare il vocabolario per la risposta.}
	
	% ============================================================
	\newpage
	\section{Varietà Simplettiche e Spazio delle Fasi}
	% ============================================================
	
	\nota{Scopo del capitolo: introdurre lo spazio delle fasi come oggetto geometrico
		intrinseco, indipendente dalla scelta di coordinate. Il teorema di Darboux mostra
		che le coordinate canoniche $(q,p)$ sono localmente universali; la struttura
		globale è la 2-forma $\omega$, che riapparirà identica come parte immaginaria del
		prodotto scalare di $\mathcal{H}$ nel Cap.\ 9. Senza questo capitolo, il collegamento
		con Ashtekar--Schilling sarebbe soltanto un'analogia; con questo, è un'identità.}
	
	% ----------------------------------------------------------
	\subsection{Geometria Simplettica}
	% ----------------------------------------------------------
	
	\fonti{Abraham \& Marsden, \textit{Foundations of Mechanics}; Cannas da Silva,
		\textit{Lectures on Symplectic Geometry}.}
	
	\begin{itemize}
		\item \textbf{Varietà simplettica:} definizione di coppia $(M, \omega)$ con $\omega$
		2-forma chiusa ($d\omega = 0$) e non degenere; prime conseguenze (dimensione
		pari, orientazione canonica).
		\item \textbf{Richiami di geometria differenziale:} campi tensoriali di tipo $(r,s)$
		su $TM$; forme alternanti; isomorfismi musicali $\flat: TM \to T^*M$ e
		$\sharp: T^*M \to TM$ indotti da $\omega$.
		\item \textbf{Campi vettoriali hamiltoniani:} dato $f \in C^\infty(M)$, definizione
		di $X_f$ tramite $\iota_{X_f}\omega = df$; esistenza e unicità garantite dalla
		non degenerazione di $\omega$.
		\item \textbf{Teorema di Darboux:} in ogni punto di $(M, \omega)$ esistono
		coordinate locali $(q^i, p_i)$ tali che $\omega = \sum_i dq^i \wedge dp_i$;
		le coordinate canoniche non sono un assioma --- sono una conseguenza locale.
	\end{itemize}
	
	% ----------------------------------------------------------
	\subsection{Simplettomorfismi e Invarianza della Struttura}
	% ----------------------------------------------------------
	
	\fonti{Abraham \& Marsden, \textit{Foundations of Mechanics}.}
	
	\begin{itemize}
		\item \textbf{Simplettomorfismo:} diffeomorfismo $\varphi: M \to M$ tale che
		$\varphi^*\omega = \omega$; gruppo $\mathrm{Symp}(M, \omega)$ delle simmetrie
		della struttura simplettica.
		\item \textbf{Teorema di Liouville:} il flusso di qualsiasi campo hamiltoniano
		è un gruppo a un parametro di simplettomorfismi; la misura di Liouville
		$\omega^n/n!$ è invariante.
		\item \textbf{Collegamento prospettico:} nella Parte~III, l'analogo quantistico
		dei simplettomorfismi sarà il gruppo unitario $U(\mathcal{H})$; il teorema di
		Stone (Cap.\ 8) stabilirà che ogni gruppo unitario fortemente continuo è
		un gruppo hamiltoniano nel senso di questo capitolo.
	\end{itemize}
	
	% ============================================================
	\newpage
	\section{Parentesi di Poisson e Algebra delle Osservabili}
	% ============================================================
	
	\nota{Scopo del capitolo: identificare le osservabili classiche come algebra di Lie,
		non solo come anello di funzioni. Questo è il punto tecnico cruciale del lato
		classico: la struttura che si vuole spiegare nel quantistico --- il Lie-bracket ---
		è già presente qui. Il primo momento del qubit (§\ref{sec:qubit1}) viene introdotto
		per fissare l'esempio che accompagnerà tutta la tesi: le funzioni coordinate su $S^2$
		con le loro parentesi di Poisson sono la controparte classica delle matrici di Pauli.}
	
	% ----------------------------------------------------------
	\subsection{$C^\infty(M)$ come Algebra di Lie}
	% ----------------------------------------------------------
	
	\fonti{Abraham \& Marsden, \textit{Foundations of Mechanics}; Cannas da Silva,
		\textit{Lectures on Symplectic Geometry}.}
	
	\begin{itemize}
		\item \textbf{Parentesi di Poisson:} definizione $\{f, g\} = \omega(X_f, X_g)$;
		proprietà: bilinearità, antisimmetria, regola di Leibniz.
		\item \textbf{Identità di Jacobi:} dimostrazione che $\{\{f,g\},h\} + \{\{g,h\},f\}
		+ \{\{h,f\},g\} = 0$ è matematicamente equivalente alla chiusura $d\omega = 0$;
		questo punto anticipa la rigidità strutturale che riapparirà su $\mathbb{P}\mathcal{H}$.
		\item \textbf{$C^\infty(M)$ come algebra di Lie:} $(C^\infty(M), \{\cdot,\cdot\})$
		soddisfa tutti gli assiomi di algebra di Lie; il Lie-bracket classico è la
		struttura astratta di cui il commutatore quantistico $\frac{1}{i\hbar}[\cdot,\cdot]$
		è l'omologo non commutativo.
		\item \textbf{Mappa $f \mapsto X_f$:} omomorfismo di algebre di Lie da
		$(C^\infty(M), \{\cdot,\cdot\})$ a $(\mathfrak{X}(M), [\cdot,\cdot])$; nucleo
		dato dalle funzioni costanti.
	\end{itemize}
	
	% ----------------------------------------------------------
	\subsection{Campi Hamiltoniani e Flusso di Poisson}
	% ----------------------------------------------------------
	
	\fonti{Abraham \& Marsden, \textit{Foundations of Mechanics}.}
	
	\begin{itemize}
		\item \textbf{Flusso di Poisson:} integrazione di $X_H$ dà una famiglia $\{\phi_t^H\}$
		di simplettomorfismi; le equazioni di Hamilton $\dot{q} = \partial_p H$,
		$\dot{p} = -\partial_q H$ sono il caso locale tramite Darboux.
		\item \textbf{Evoluzione di un'osservabile:} $\frac{d}{dt}(f \circ \phi_t^H) =
		\{f, H\} \circ \phi_t^H$; la dinamica classica è interamente codificata nel
		Lie-bracket.
	\end{itemize}
	
	% ----------------------------------------------------------
	\subsection{Qubit I: $S^2$ come Spazio delle Fasi dello Spin Classico}
	\label{sec:qubit1}
	% ----------------------------------------------------------
	
	\fonti{Abraham \& Marsden, \textit{Foundations of Mechanics}; Perelomov,
		\textit{Generalized Coherent States}.}
	
	\begin{itemize}
		\item \textbf{Spazio delle fasi dello spin classico:} per un sistema con simmetria
		$SU(2)$, lo spazio delle fasi ridotto è la 2-sfera $S^2$ con forma simplettica
		$\omega = \sin\theta\, d\theta \wedge d\phi$ (forma d'area normalizzata).
		\item \textbf{Parentesi di Poisson su $S^2$:} in coordinate $(x_1, x_2, x_3)$
		con il vincolo $\sum x_i^2 = 1$, si ha $\{x_i, x_j\} = \varepsilon_{ijk} x_k$;
		queste sono le parentesi di Poisson dell'algebra di Lie $\mathfrak{su}(2)$.
		\item \textbf{Anticipazione:} nel Cap.\ 7 si mostrerà che il commutatore
		$[\sigma_i, \sigma_j] = 2i\varepsilon_{ijk}\sigma_k$ su $M_2(\mathbb{C})$ è la
		stessa algebra di Lie --- non per analogia, ma perché entrambe le algebre
		($C^\infty(S^2)$ e $M_2(\mathbb{C})$) sono rappresentazioni distinte di
		$\mathfrak{su}(2)$.
	\end{itemize}
	
	% ============================================================
	\newpage
	\section{Flusso Hamiltoniano e Teorema di Noether}
	% ============================================================
	
	\nota{Scopo del capitolo: formulare il teorema di Noether in forma invariante, tramite
		la mappa momento. Questa formulazione --- conservazione di $\mu$ lungo il flusso di
		$X_H$ quando $H$ è $G$-invariante --- è identica nella meccanica razionale e in quella
		quantistica. Il Cap.\ 10 riprenderà questa struttura e mostrerà che $[H, A] = 0$
		implica $\frac{d}{dt}\langle A \rangle = 0$ per esattamente lo stesso motivo geometrico.
		La formulazione coordinata-libera qui è necessaria perché solo quella si trasferisce
		direttamente.}
	
	% ----------------------------------------------------------
	\subsection{Gruppi di Lie e Azioni Simplettiche}
	% ----------------------------------------------------------
	
	\fonti{Abraham \& Marsden, \textit{Foundations of Mechanics}; Moretti,
		\textit{Spectral Theory and Quantum Mechanics}.}
	
	\begin{itemize}
		\item \textbf{Gruppo di Lie $G$:} definizione; algebra di Lie $\mathfrak{g} = T_eG$
		con parentesi di Lie; mappa esponenziale $\exp: \mathfrak{g} \to G$; esempi
		rilevanti $U(1)$, $SU(2)$, $U(n)$.
		\item \textbf{Azione simplettica:} $\Phi: G \times M \to M$ liscia con
		$\Phi_g^* \omega = \omega$ per ogni $g \in G$; campo vettoriale generatore
		infinitesimale $\xi_M = \frac{d}{dt}\big|_{t=0} \Phi_{\exp(t\xi)}$.
		\item \textbf{Sottogruppi a un parametro:} corrispondenza biunivoca con elementi
		di $\mathfrak{g}$; integrazione come flusso del campo generatore.
	\end{itemize}
	
	% ----------------------------------------------------------
	\subsection{Mappa Momento e Teorema di Noether}
	% ----------------------------------------------------------
	
	\fonti{Abraham \& Marsden, \textit{Foundations of Mechanics}.}
	
	\begin{itemize}
		\item \textbf{Mappa momento:} applicazione $\mu: M \to \mathfrak{g}^*$ definita da
		$d\langle\mu, \xi\rangle = \iota_{\xi_M}\omega$ per ogni $\xi \in \mathfrak{g}$;
		$\langle\mu, \xi\rangle$ è la funzione hamiltoniana del campo $\xi_M$.
		\item \textbf{Teorema di Noether geometrico:} se $H$ è $G$-invariante (cioè
		$\{H, \langle\mu, \xi\rangle\} = 0$ per ogni $\xi$), allora $\mu$ è conservata
		lungo il flusso di $X_H$:
		\[
		\frac{d}{dt}\langle\mu \circ \phi_t^H, \xi\rangle = \{\langle\mu,\xi\rangle, H\} = 0.
		\]
		\item \textbf{Formulazione invariante:} simmetria $\Leftrightarrow$ conservazione,
		mediata dal Lie-bracket; questa formulazione non dipende da coordinate né dalla
		dimensione dello spazio --- si applica identicamente a $M$ finito-dimensionale e
		a $\mathbb{P}\mathcal{H}$ infinito-dimensionale.
	\end{itemize}
	
	% ============================================================
	\newpage
	\part*{Parte II --- Il Quadro Algebrico}
	\addcontentsline{toc}{part}{Parte II: Il Quadro Algebrico}
	% ============================================================
	
	\nota{La Parte~II stabilisce il linguaggio algebrico che unifica le due teorie.
		I tre capitoli hanno una dipendenza logica stretta: il Cap.\ 4 mostra che entrambe
		le teorie sono $C^*$-algebre; il Cap.\ 5 mostra perché le CCR canoniche non
		vivono in una $C^*$-algebra e come l'algebra di Weyl risolve il problema;
		il Cap.\ 6 usa l'algebra di Weyl del Cap.\ 5 come esempio principale della
		strict deformation quantization di Rieffel. Saltare uno qualsiasi di questi
		capitoli lascia un buco logico.}
	
	% ============================================================
	\newpage
	\section{$C^*$-Algebre: Struttura Comune di MR e MQ}
	% ============================================================
	
	\nota{Questo è il capitolo cardine della tesi. Senza di esso i capitoli successivi
		non hanno fondamento. L'obiettivo è mostrare che la meccanica razionale e la meccanica
		quantistica sono entrambe $C^*$-algebre --- non per scelta, ma per classificazione
		(teorema di Gelfand nel caso commutativo, teorema di Gelfand--Naimark nel caso
		generale). Lo spazio di Hilbert $\mathcal{H}$ emerge dal teorema GNS come spazio
		di rappresentazione, non come assioma fisico: questa è la prima risposta parziale
		alla domanda della tesi.}
	
	% ----------------------------------------------------------
	\subsection{Definizione Assiomatica}
	% ----------------------------------------------------------
	
	\fonti{Bratteli \& Robinson, \textit{Operator Algebras and Quantum Statistical
			Mechanics}, Vol.~I; Strocchi, \textit{An Introduction to the Mathematical
			Structure of Quantum Mechanics}.}
	
	\begin{itemize}
		\item \textbf{$C^*$-algebra:} algebra di Banach $\mathcal{A}$ con involuzione $*$
		soddisfacente la condizione C*: $\|a^*a\| = \|a\|^2$; esempi fondamentali
		$B(\mathcal{H})$, $C_0(X)$, $M_n(\mathbb{C})$.
		\item \textbf{Elementi autoaggiunti:} $a = a^*$ implica $\sigma(a) \subset \mathbb{R}$;
		ruolo fisico come osservabili con valori reali.
		\item \textbf{Norma determinata dalla struttura algebrica:} l'identità
		$\|a\|^2 = r(a^*a)$ (raggio spettrale) mostra che la norma non è una scelta
		libera --- è determinata dai soli dati algebrici. Questo è un primo segnale
		del perché le due teorie ``si assomigliano'': usano la stessa struttura.
		\item \textbf{Stato come funzionale positivo normalizzato:} $\omega: \mathcal{A}
		\to \mathbb{C}$ lineare con $\omega(a^*a) \geq 0$ e $\omega(I) = 1$; il valore
		di aspettazione $\omega(a)$ come dato fisico primitivo (approccio Strocchi:
		nessuno spazio di Hilbert assunto a priori).
	\end{itemize}
	
	% ----------------------------------------------------------
	\subsection{MR come $C^*$-Algebra Commutativa: Teorema di Gelfand}
	% ----------------------------------------------------------
	
	\fonti{Bratteli \& Robinson, Vol.~I; Strocchi, \textit{An Introduction to the
			Mathematical Structure of Quantum Mechanics}.}
	
	\begin{itemize}
		\item \textbf{$C_0(M)$ come $C^*$-algebra:} le funzioni continue che si
		annullano all'infinito su una varietà simplettica $M$, con norma sup e
		involuzione data dalla coniugazione complessa, formano una $C^*$-algebra
		commutativa.
		\item \textbf{Teorema di Gelfand:} ogni $C^*$-algebra \emph{commutativa} è
		isometricamente isomorfa a $C_0(X)$ per qualche spazio localmente compatto
		di Hausdorff $X$; la MR non ``usa'' $C^*$-algebre --- la MR \emph{è} la
		teoria delle $C^*$-algebre commutative, per classificazione.
		\item \textbf{Osservabili classiche come elementi autoaggiunti di $C_0(M)$:}
		le funzioni reali $f \in C_0(M)$ sono gli elementi autoaggiunti di questa
		$C^*$-algebra; il prodotto puntuale è il prodotto associativo (commutativo).
	\end{itemize}
	
	% ----------------------------------------------------------
	\subsection{MQ come $C^*$-Algebra Non Commutativa: Teorema di Gelfand--Naimark}
	% ----------------------------------------------------------
	
	\fonti{Bratteli \& Robinson, Vol.~I; Reed \& Simon, \textit{Methods of Modern
			Mathematical Physics}, Vol.~I.}
	
	\begin{itemize}
		\item \textbf{$B(\mathcal{H})$ come $C^*$-algebra:} gli operatori limitati su
		$\mathcal{H}$ con norma operatoriale e involuzione data dall'aggiunto formano
		una $C^*$-algebra non commutativa.
		\item \textbf{Osservabili quantistiche come elementi autoaggiunti:} $\hat{A} =
		\hat{A}^*$; dinamica come automorfismo dell'algebra (preserva la struttura).
		\item \textbf{Teorema di Gelfand--Naimark:} ogni $C^*$-algebra astratta si
		rappresenta isometricamente come sottoalgebra di $B(\mathcal{H})$ per qualche
		$\mathcal{H}$; lo spazio di Hilbert non è un assioma fisico --- è il supporto
		naturale di qualsiasi $C^*$-algebra.
	\end{itemize}
	
	% ----------------------------------------------------------
	\subsection{Teorema GNS: lo Spazio di Hilbert come Conseguenza}
	% ----------------------------------------------------------
	
	\fonti{Reed \& Simon, Vol.~I; Moretti, \textit{Spectral Theory and Quantum Mechanics};
		Strocchi, \textit{An Introduction to the Mathematical Structure of Quantum Mechanics}.}
	
	\begin{itemize}
		\item \textbf{Prodotto interno su $\mathcal{A}$:} dato uno stato $\omega$,
		$\langle a, b\rangle_\omega = \omega(a^*b)$; la positività di $\omega$ garantisce
		la proprietà di prodotto interno (salvo quoziente per l'ideale di Gelfand
		$\mathcal{I}_\omega$).
		\item \textbf{Spazio di Hilbert GNS:} $\mathcal{H}_\omega$ come completamento di
		$\mathcal{A}/\mathcal{I}_\omega$; rappresentazione $\pi_\omega: \mathcal{A} \to
		B(\mathcal{H}_\omega)$ per moltiplicazione a sinistra; vettore ciclico $\Omega_\omega
		= [I]$.
		\item \textbf{Teorema GNS:} ogni stato $\omega$ su una $C^*$-algebra $\mathcal{A}$
		determina (a meno di equivalenza unitaria) un'unica rappresentazione
		$(\mathcal{H}_\omega, \pi_\omega, \Omega_\omega)$ con
		$\omega(a) = \langle\Omega_\omega, \pi_\omega(a)\Omega_\omega\rangle$;
		$\mathcal{H}$ \emph{emerge} dalla struttura algebrica, non è assunto.
		\item \textbf{Conclusione del capitolo:} le due prime domande della tesi hanno
		già una risposta parziale: MR e MQ sono due $C^*$-algebre, una commutativa e
		una no; il Lie-bracket $[a,b] = ab - ba$ è struttura interna di ogni
		$C^*$-algebra non commutativa. Il Cap.\ 5 precisere perché questo Lie-bracket
		non è però sufficiente a realizzare le CCR $[q,p] = i\hbar I$.
	\end{itemize}
	
	% ----------------------------------------------------------
	\subsection{Qubit II: $M_2(\mathbb{C})$ come $C^*$-Algebra}
	% ----------------------------------------------------------
	
	\fonti{Bratteli \& Robinson, Vol.~I; Landsman, \textit{Mathematical Topics Between
			Classical and Quantum Mechanics}.}
	
	\begin{itemize}
		\item \textbf{Algebra degli osservabili:} $\mathcal{A} = M_2(\mathbb{C})$; matrici
		hermitiane come osservabili; non-commutatività esplicita:
		$[\sigma_x, \sigma_y] = 2i\sigma_z$.
		\item \textbf{Stati come matrici densità:} $\rho \in M_2(\mathbb{C})$ con
		$\rho \geq 0$, $\mathrm{tr}(\rho) = 1$; stati puri $\rho = |\psi\rangle\langle\psi|$
		con $\rho^2 = \rho$; stati misti con $\mathrm{tr}(\rho^2) < 1$.
		\item \textbf{GNS esplicito:} per lo stato traciale $\omega(a) = \frac{1}{2}\mathrm{tr}(a)$,
		lo spazio GNS è $\mathbb{C}^2$ con la rappresentazione standard; il vettore
		ciclico è $\Omega = \frac{1}{\sqrt{2}}(1, 1)^T$. Calcolo esplicito della
		costruzione.
	\end{itemize}
	
	% ============================================================
	\newpage
	\section{Wintner--Wielandt e Algebra di Weyl}
	% ============================================================
	
	\nota{Scopo del capitolo: mostrare perché le relazioni di commutazione canoniche
		(CCR) $[q,p] = i\hbar I$ non possono essere realizzate in una $C^*$-algebra
		(teorema di Wintner--Wielandt), e come l'algebra di Weyl risolve il problema
		passando agli unitari esponenziali. Il teorema di Stone--von Neumann, che fornisce
		l'unicità della rappresentazione, viene enunciato qui ma \emph{usato} nel Cap.\ 8,
		dove serve realmente. La sezione finale (§\ref{sec:liebracketinternal}) è il
		collegamento diretto verso il Cap.\ 7: il Lie-bracket $[a,b] = ab - ba$ è struttura
		interna di ogni $C^*$-algebra non commutativa --- non è un postulato aggiunto.}
	
	% ----------------------------------------------------------
	\subsection{Il Problema delle CCR}
	% ----------------------------------------------------------
	
	\fonti{Strocchi, \textit{An Introduction to the Mathematical Structure of Quantum
			Mechanics}; Bratteli \& Robinson, Vol.~II.}
	
	\begin{itemize}
		\item \textbf{Inquadramento:} si cerca una $C^*$-algebra contenente $q, p$ con
		$[q, p] = i\hbar I$; se tali operatori fossero limitati, si potrebbe prendere
		la traccia: $\mathrm{tr}([q,p]) = 0 \neq \mathrm{tr}(i\hbar I)$. Contraddizione.
		\item \textbf{Necessità di uscire dagli operatori limitati:} le CCR richiedono
		operatori illimitati; ma gli operatori illimitati non formano una $C^*$-algebra.
		Il problema è strutturale, non tecnico.
	\end{itemize}
	
	% ----------------------------------------------------------
	\subsection{Teorema di Wintner--Wielandt}
	% ----------------------------------------------------------
	
	\fonti{Strocchi, \textit{An Introduction to the Mathematical Structure of Quantum
			Mechanics}; Moretti, \textit{Spectral Theory and Quantum Mechanics}, p.~605.}
	
	\begin{itemize}
		\item \textbf{Enunciato:} in uno spazio di Banach non esistono operatori
		\emph{limitati} $A, B$ tali che $AB - BA = I$; dimostrazione per induzione
		sulla traccia di $B^n A - AB^n$.
		\item \textbf{Conseguenza (\textit{momento drammatico}):} le CCR $[q,p] = i\hbar I$
		non possono essere soddisfatte in nessuna $C^*$-algebra unitale. La
		non-commutatività quantistica, nella forma delle CCR, è incompatibile con la
		struttura di algebra di Banach. Bisogna riformulare il problema.
		\item \textbf{Wintner--Wielandt non proibisce la quantizzazione:} proibisce
		che le CCR siano realizzate da operatori limitati; l'algebra di Weyl aggira
		questo ostacolo passando alle CCR in forma esponenziale.
	\end{itemize}
	
	% ----------------------------------------------------------
	\subsection{Algebra di Weyl come Soluzione}
	% ----------------------------------------------------------
	
	\fonti{Bratteli \& Robinson, Vol.~II; Strocchi, \textit{An Introduction to the
			Mathematical Structure of Quantum Mechanics}.}
	
	\begin{itemize}
		\item \textbf{Operatori unitari di Weyl:} $W(\alpha) = e^{i(s\hat{q} + t\hat{p})}$
		per $\alpha = (s, t) \in \mathbb{R}^2$; sono limitati (unitari) anche se $\hat{q}$
		e $\hat{p}$ non lo sono.
		\item \textbf{Relazioni di Weyl:} $W(\alpha)W(\beta) = e^{\frac{i}{2}\omega(\alpha,\beta)}
		W(\alpha + \beta)$ con $\omega(\alpha,\beta) = st' - ts'$; queste sono le CCR
		in forma esponenziale, evitando il problema di dominio degli operatori illimitati.
		\item \textbf{$C^*$-algebra di Weyl $\mathrm{CCR}(\mathbb{R}^2)$:} completamento
		in norma C* dell'algebra involutiva generata dai $W(\alpha)$; la norma è unica
		(la condizione C* impone $\|W(\alpha)\| = 1$).
		\item \textbf{Recupero delle CCR di Heisenberg:} tramite il teorema di Stone,
		i generatori infinitesimali degli unitari $W(s, 0)$ e $W(0, t)$ soddisfano
		$[q, p] = i\hbar I$ nel senso degli operatori non limitati sui rispettivi domini.
		\item \textbf{Teorema di Stone--von Neumann} (enunciato; usato nel Cap.~8): ogni
		rappresentazione irriducibile fortemente continua di $\mathrm{CCR}(\mathbb{R}^{2n})$
		è unitariamente equivalente alla rappresentazione di Schrödinger su
		$L^2(\mathbb{R}^n)$; unicità cinematica della quantizzazione canonica per sistemi
		a finiti gradi di libertà.
	\end{itemize}
	
	% ----------------------------------------------------------
	\subsection{Il Lie-Bracket come Struttura Interna}
	\label{sec:liebracketinternal}
	% ----------------------------------------------------------
	
	\fonti{Bratteli \& Robinson, Vol.~I.}
	
	\begin{itemize}
		\item \textbf{$[a,b] = ab - ba$ in ogni $C^*$-algebra:} il commutatore di due
		elementi è antisimmetrico e soddisfa l'identità di Jacobi per costruzione
		algebrica; ogni $C^*$-algebra non commutativa è automaticamente un'algebra
		di Lie rispetto a $[\cdot, \cdot]$.
		\item \textbf{Conseguenza fondamentale:} il Lie-bracket quantistico non è un
		assioma aggiunto alla $C^*$-algebra --- è struttura interna. Questo prepara
		la risposta del Cap.\ 7: la corrispondenza tra commutatore e parentesi di Poisson
		non è un'analogia imposta, ma una conseguenza del fatto che entrambi sono
		Lie-bracket della rispettiva $C^*$-algebra.
	\end{itemize}
	
	% ============================================================
	\newpage
	\section{Strict Deformation Quantization (Rieffel)}
	% ============================================================
	
	\nota{Scopo del capitolo: costruire il ponte continuo tra classico e quantistico.
		Il capitolo risponde alla domanda: in che senso $A_0$ e $A_\hbar$ sono ``lo stesso
		oggetto''? La struttura interna è: (1) problema --- la deformazione formale non basta;
		(2) definizione --- Rieffel fornisce la risposta corretta tramite $C^*$-algebre vere;
		(3) esempio --- l'algebra di Weyl del Cap.\ 5 è l'istanza principale; il qubit
		è il caso finito-dimensionale calcolabile esplicitamente. La condizione di Dirac
		è il cuore tecnico: è essa che garantisce la persistenza delle strutture classiche
		nel senso normato, non per analogia.}
	
	% ----------------------------------------------------------
	\subsection{Introduzione: il Parametro $\hbar$}
	% ----------------------------------------------------------
	
	\fonti{Landsman, \textit{Mathematical Topics Between Classical and Quantum Mechanics}.}
	
	\begin{itemize}
		\item \textbf{(i) $\hbar$ come costante fisica fissata:} presupposto per la
		cinematica degli operatori e la validità delle CCR (Capp.\ 4--5).
		\item \textbf{(ii) $\hbar$ come parametro formale:} trattazione di $\hbar$ come
		variabile nelle serie di potenze formali $\mathbb{R}[[\hbar]]$, base della
		deformazione formale (§\ref{sec:formalDQ}).
		\item \textbf{(iii) $\hbar$ come indice continuo:} $\hbar \in [0, 1]$ parametrizza
		le fibre di un campo continuo di $C^*$-algebre; il limite classico $\hbar \to 0$
		è un limite normato nella topologia C* (strict deformation quantization).
		\item \textbf{Interpretazione del limite classico:} il rigore matematico appartiene
		solo al regime (iii); il ``taglio di Heisenberg'' del Cap.\ 10 è il regime in cui
		$\|Q_\hbar(f) - Q_0(f)\|$ diventa trascurabile nella norma C*.
	\end{itemize}
	
	% ----------------------------------------------------------
	\subsection{Perché la Deformazione Formale Non Basta}
	\label{sec:formalDQ}
	% ----------------------------------------------------------
	
	\fonti{Landsman, \textit{Mathematical Topics Between Classical and Quantum Mechanics};
		Kontsevich, \textit{Deformation Quantization of Poisson Manifolds}
		(Lett.\ Math.\ Phys.\ 66, 2003).}
	
	\begin{itemize}
		\item \textbf{Star-product formale (BFFLS):} mappa bilineare $\star_\hbar:
		C^\infty(M)[[\hbar]] \times C^\infty(M)[[\hbar]] \to C^\infty(M)[[\hbar]]$ con
		$f \star_\hbar g = \sum_{k=0}^\infty \hbar^k B_k(f,g)$; $B_0(f,g) = fg$;
		$B_1(f,g) - B_1(g,f) = i\{f,g\}$.
		\item \textbf{Problema:} le serie formali in $\hbar$ non definiscono una
		$C^*$-algebra; non esiste in generale una norma rispetto a cui la serie converga;
		il limite $\hbar \to 0$ non ha senso normato. La deformazione formale è uno
		strumento algebrico, non un oggetto fisico.
		\item \textbf{Necessità della strict DQ:} si vuole che per ogni valore di $\hbar$
		si abbia una \emph{vera} $C^*$-algebra, e che la famiglia sia continua in $\hbar$
		nella topologia della norma. Questo è ciò che Rieffel costruisce.
	\end{itemize}
	
	% ----------------------------------------------------------
	\subsection{Definizione di Rieffel e Condizione di Dirac}
	% ----------------------------------------------------------
	
	\fonti{Rieffel, \textit{Deformation Quantization for Actions of $\mathbb{R}^d$}
		(Mem.\ AMS, 1993); Landsman, \textit{Mathematical Topics Between Classical and
			Quantum Mechanics}.}
	
	\begin{itemize}
		\item \textbf{Definizione (Rieffel, 1993):} una strict deformation quantization
		di $(A_0, \{\cdot,\cdot\})$ è una famiglia $\{A_\hbar\}_{\hbar \in I}$ di
		$C^*$-algebre (con $A_0 = C_0(M)$) e mappe $Q_\hbar: \widetilde{A}_0 \to A_\hbar$
		definite su un'algebra densa $\widetilde{A}_0 \subset A_0$, tali che:
		\begin{itemize}
			\item[(i)] $Q_\hbar(\bar{f}) = Q_\hbar(f)^*$ (compatibilità con l'involuzione);
			\item[(ii)] $\|Q_\hbar(f)\|_{A_\hbar} \to \|f\|_{A_0}$ per $\hbar \to 0$
			(fedeltà nel limite classico);
			\item[(iii)] \emph{Condizione di Dirac:}
			$\bigl\|\tfrac{1}{i\hbar}[Q_\hbar(f), Q_\hbar(g)] - Q_\hbar(\{f,g\})\bigr\|_{A_\hbar}
			\to 0$ per $\hbar \to 0$.
		\end{itemize}
		\item \textbf{Cosa garantisce la condizione di Dirac:} non che il commutatore
		sia \emph{uguale} alla parentesi di Poisson, ma che la loro differenza sia piccola
		in norma C* per $\hbar$ piccolo. Le strutture classiche persistono nel quantistico
		come \emph{limite normato}, non come uguaglianza.
		\item \textbf{Campo continuo di $C^*$-algebre (Landsman):} la famiglia
		$\{A_\hbar\}$ si organizza in un campo continuo; classicità e quantisticità non
		sono separate da una discontinuità ontologica --- il confine è un limite
		asintotico nella topologia della norma.
	\end{itemize}
	
	% ----------------------------------------------------------
	\subsection{Esempio: Quantizzazione di Weyl su $\mathbb{R}^2$}
	% ----------------------------------------------------------
	
	\fonti{Rieffel, \textit{Deformation Quantization for Actions of $\mathbb{R}^d$};
		Landsman, \textit{Mathematical Topics Between Classical and Quantum Mechanics}.}
	
	\begin{itemize}
		\item \textbf{Setup:} $A_0 = C_0(\mathbb{R}^2)$; $A_\hbar = \mathrm{CCR}(\mathbb{R}^2)$
		(algebra di Weyl del Cap.\ 5); $Q_\hbar$ = quantizzazione di Weyl
		$Q_\hbar^W(f) = \frac{1}{(2\pi\hbar)}\int f(q,p) W(q,p)\, dq\, dp$.
		\item \textbf{Verifica delle tre condizioni di Rieffel:}
		\begin{itemize}
			\item[(i)] compatibilità con l'involuzione: $Q_\hbar^W(\bar{f}) = Q_\hbar^W(f)^*$
			per le proprietà degli operatori di Weyl;
			\item[(ii)] limite della norma: $\|Q_\hbar^W(f)\| \to \|f\|_\infty$ per
			$\hbar \to 0$, verificato tramite la stima di Calderón--Vaillancourt;
			\item[(iii)] condizione di Dirac: segue dal calcolo esplicito del commutatore
			degli operatori di Weyl.
		\end{itemize}
		\item \textbf{Conclusione:} la quantizzazione di Weyl è una strict deformation
		quantization di $(\mathbb{R}^2, \omega_{\mathrm{can}})$ nel senso di Rieffel.
	\end{itemize}
	
	% ----------------------------------------------------------
	\subsection{Qubit III: Famiglia $A_\hbar$ che Interpola $C(S^2) \to M_2(\mathbb{C})$}
	% ----------------------------------------------------------
	
	\fonti{Landsman, \textit{Mathematical Topics Between Classical and Quantum Mechanics};
		Rieffel, \textit{Deformation Quantization for Actions of $\mathbb{R}^d$}.}
	
	\begin{itemize}
		\item \textbf{Setup finito-dimensionale:} $A_0 = C(S^2)$;
		$A_{\hbar=1} = M_2(\mathbb{C})$; la famiglia interpola tra la $C^*$-algebra
		classica dello spin e quella quantistica.
		\item \textbf{Verifica esplicita delle condizioni di Rieffel:} nel caso
		finito-dimensionale tutte le norme sono equivalenti; le tre condizioni si
		verificano con calcoli elementari.
		\item \textbf{Significato fisico:} per $\hbar = 1$ si ha il qubit; per $\hbar \to 0$
		si recupera l'algebra delle funzioni sullo spin classico. Le strutture di
		Poisson su $S^2$ e il Lie-bracket di $M_2(\mathbb{C})$ sono connessi da questa
		famiglia continua, non solo per analogia.
	\end{itemize}
	
	% ============================================================
	\newpage
	\part*{Parte III --- Le Risposte}
	\addcontentsline{toc}{part}{Parte III: Le Risposte}
	% ============================================================
	
	\nota{La Parte~III risponde alla domanda della tesi capitolo per capitolo.
		Ogni capitolo usa tutto il lavoro delle Parti~I e~II per rispondere a una delle
		domande implicite: perché il Lie-bracket? (Cap.\ 7); perché il flusso hamiltoniano?
		(Cap.\ 8); perché la struttura simplettica? (Cap.\ 9); perché Noether?
		e che ne è del taglio di Heisenberg? (Cap.\ 10). La risposta sintetica è nel
		Cap.\ 10: le strutture non riemergono perché non se ne sono mai andate --- sono
		struttura dell'algebra, non della fisica specifica.}
	
	% ============================================================
	\newpage
	\section{La Corrispondenza Classico-Quantistica: Tre Livelli}
	% ============================================================
	
	\nota{Scopo del capitolo: rispondere alla domanda sul Lie-bracket. Il capitolo è
		strutturato in tre livelli di profondità crescente. Il livello~1 usa la condizione
		di Dirac di Rieffel: la corrispondenza tra commutatore e parentesi di Poisson non
		è un'analogia scoperta a posteriori, è una proprietà richiesta dalla definizione
		di strict DQ. Il livello~2 (Groenewold--van Hove) delimita la portata di questa
		risposta: la corrispondenza è asintotica, non esatta. Il livello~3 è la risposta
		più profonda, indipendente da Rieffel: il Lie-bracket è struttura interna di ogni
		$C^*$-algebra non commutativa (§\ref{sec:liebracketinternal}), non un assioma aggiunto.}
	
	% ----------------------------------------------------------
	\subsection{Livello 1 --- Corrispondenza come Continuità}
	% ----------------------------------------------------------
	
	\fonti{Rieffel, \textit{Deformation Quantization for Actions of $\mathbb{R}^d$};
		Landsman, \textit{Mathematical Topics Between Classical and Quantum Mechanics}.}
	
	\begin{itemize}
		\item \textbf{La condizione di Dirac è un teorema, non un principio:} per ogni
		strict DQ nel senso di Rieffel, la condizione (iii) del Cap.\ 6 afferma
		\emph{per definizione} che
		\[
		\Bigl\|\frac{1}{i\hbar}[Q_\hbar(f), Q_\hbar(g)] - Q_\hbar(\{f,g\})\Bigr\|
		\xrightarrow{\hbar \to 0} 0.
		\]
		La corrispondenza di Dirac non è un principio imposto alla teoria:
		è una proprietà richiesta dalla definizione stessa di quantizzazione, e verificata
		esplicitamente nel caso di Weyl.
		\item \textbf{Senso preciso della corrispondenza:} il commutatore $[Q_\hbar(f),
		Q_\hbar(g)]/i\hbar$ \emph{converge in norma} a $Q_\hbar(\{f,g\})$ per
		$\hbar \to 0$. Non è un'analogia né un'approssimazione informale: è una
		disuguaglianza normata.
	\end{itemize}
	
	% ----------------------------------------------------------
	\subsection{Livello 2 --- Ostacolo di Groenewold--van Hove}
	% ----------------------------------------------------------
	
	\fonti{Moretti, \textit{Spectral Theory and Quantum Mechanics}, p.~605;
		Gotay, \textit{Obstructions to Quantization} (arXiv:math-ph/9809011).}
	
	\begin{itemize}
		\item \textbf{Teorema di Groenewold--van Hove:} non esiste nessuna mappa lineare
		$Q: \mathcal{P}(\mathbb{R}^{2n}) \to L(\mathcal{H})$ sull'algebra dei polinomi
		che soddisfi simultaneamente le regole di Dirac (D1)--(D4); il conflitto emerge
		già sui monomi di grado $\leq 4$.
		\item \textbf{Compatibilità con Rieffel:} la condizione di Dirac di Rieffel è
		\emph{asintotica} (limite $\hbar \to 0$), non un'uguaglianza esatta per $\hbar > 0$.
		Groenewold--van Hove proibisce l'uguaglianza esatta su tutto $C^\infty$; non
		proibisce la continuità asintotica.
		\item \textbf{Portata della risposta:} la corrispondenza è precisa e normata
		nel limite semiclassico; per $\hbar > 0$ fissato, esistono osservabili per cui
		la quantizzazione non può essere un omomorfismo esatto di algebre di Lie.
	\end{itemize}
	
	% ----------------------------------------------------------
	\subsection{Livello 3 --- Il Lie-Bracket è Intrinseco}
	% ----------------------------------------------------------
	
	\fonti{Bratteli \& Robinson, Vol.~I.}
	
	\begin{itemize}
		\item \textbf{Struttura indipendente da Rieffel:} in ogni $C^*$-algebra non
		commutativa, $[a,b] = ab - ba$ soddisfa antisimmetria e identità di Jacobi
		per costruzione algebrica pura (§\ref{sec:liebracketinternal}). La corrispondenza
		con il Lie-bracket classico non dipende dal limite $\hbar \to 0$: dipende dal
		fatto che entrambe le algebre (commutativa e non) portano un Lie-bracket come
		struttura interna.
		\item \textbf{Risposta definitiva:} il commutatore quantistico non ``corrisponde''
		alla parentesi di Poisson --- \emph{sono} la stessa struttura (Lie-bracket)
		realizzata in due algebre diverse: una commutativa (la MR), una non commutativa
		(la MQ).
	\end{itemize}
	
	% ----------------------------------------------------------
	\subsection{Qubit IV: Calcolo Esplicito}
	% ----------------------------------------------------------
	
	\fonti{Bratteli \& Robinson, Vol.~I; Landsman, \textit{Mathematical Topics Between
			Classical and Quantum Mechanics}.}
	
	\begin{itemize}
		\item \textbf{Lato quantistico:} $[\sigma_i, \sigma_j] = 2i\varepsilon_{ijk}\sigma_k$
		in $M_2(\mathbb{C})$; Lie-bracket di $\mathfrak{su}(2)$ realizzato come commutatore.
		\item \textbf{Lato classico:} $\{x_i, x_j\} = \varepsilon_{ijk} x_k$ su $S^2$;
		stesso Lie-bracket di $\mathfrak{su}(2)$, realizzato come parentesi di Poisson.
		\item \textbf{Conclusione:} \emph{stessa algebra di Lie, due rappresentazioni}.
		$C^\infty(S^2)$ con $\{\cdot,\cdot\}$ e $M_2(\mathbb{C})$ con
		$\frac{1}{2i}[\cdot,\cdot]$ sono due rappresentazioni di $\mathfrak{su}(2)$;
		la condizione di Dirac di Rieffel garantisce che la seconda converge alla prima
		per $\hbar \to 0$.
	\end{itemize}
	
	% ============================================================
	\newpage
	\section{Evoluzione Unitaria e Flusso Hamiltoniano}
	% ============================================================
	
	\nota{Scopo del capitolo: rispondere alla domanda sul flusso hamiltoniano. Il teorema
		di Stone--von Neumann (enunciato nel Cap.\ 5, usato qui) garantisce che la
		rappresentazione di Schrödinger è unica. Il teorema di Stone poi mostra che ogni
		gruppo unitario fortemente continuo è un flusso hamiltoniano, nel senso del Cap.\ 3.
		L'equazione di Schrödinger non è un assioma aggiunto: è il caso quantistico del
		principio variazionale hamiltoniano. Il qubit (precessione di Larmor) rende
		esplicito tutto questo in quattro numeri complessi.}
	
	% ----------------------------------------------------------
	\subsection{Stone--von Neumann e Unicità della Rappresentazione}
	% ----------------------------------------------------------
	
	\fonti{Reed \& Simon, Vol.~I; Moretti, \textit{Spectral Theory and Quantum Mechanics}.}
	
	\begin{itemize}
		\item \textbf{Teorema di Stone--von Neumann} (da Cap.\ 5): ogni rappresentazione
		irriducibile fortemente continua di $\mathrm{CCR}(\mathbb{R}^{2n})$ è
		unitariamente equivalente alla rappresentazione di Schrödinger su $L^2(\mathbb{R}^n)$;
		$\hat{q}_j \psi = x_j \psi$, $\hat{p}_j \psi = -i\hbar\partial_{x_j}\psi$.
		\item \textbf{Perché è usato qui:} Stone--von Neumann identifica \emph{lo} spazio
		di Hilbert della quantizzazione canonica; questa identificazione è necessaria per
		applicare il teorema di Stone all'hamiltoniana specifica.
		\item \textbf{Limite di validità:} Stone--von Neumann fallisce in dimensione
		infinita (QFT); esistono rappresentazioni non equivalenti. L'approccio di questo
		capitolo è valido per sistemi a finiti gradi di libertà.
	\end{itemize}
	
	% ----------------------------------------------------------
	\subsection{Teorema di Stone: Equazione di Schrödinger come Flusso}
	% ----------------------------------------------------------
	
	\fonti{Reed \& Simon, Vol.~I; Moretti, \textit{Spectral Theory and Quantum Mechanics}.}
	
	\begin{itemize}
		\item \textbf{Teorema di Stone:} ogni gruppo unitario fortemente continuo
		$\{U(t)\}_{t \in \mathbb{R}}$ su $\mathcal{H}$ ammette un unico generatore
		autoaggiunto $H$ tale che $U(t) = e^{-iHt/\hbar}$; conversamente, ogni operatore
		autoaggiunto genera tale gruppo.
		\item \textbf{Equazione di Schrödinger come forma differenziale:}
		$i\hbar\partial_t|\psi\rangle = H|\psi\rangle$ è la derivata in $t=0$ di
		$U(t)|\psi_0\rangle$; non è un postulato indipendente --- è il teorema di Stone
		applicato al gruppo di evoluzione temporale.
		\item \textbf{Parallelo con il Cap.\ 3:} nel Cap.\ 3 ogni campo hamiltoniano
		$X_H$ genera un flusso simplettico; qui ogni operatore autoaggiunto $H$ genera
		un flusso unitario. Il teorema di Stone è l'analogo quantistico del teorema
		di esistenza e unicità del flusso hamiltoniano --- entrambi discendono dalla
		struttura di algebra di Lie della rispettiva $C^*$-algebra.
	\end{itemize}
	
	% ----------------------------------------------------------
	\subsection{Qubit V: Precessione di Larmor}
	% ----------------------------------------------------------
	
	\fonti{Landsman, \textit{Mathematical Topics Between Classical and Quantum Mechanics}.}
	
	\begin{itemize}
		\item \textbf{Hamiltoniana del qubit:} $H = -\gamma \vec{B} \cdot \vec{\sigma}/2$
		con $\vec{B}$ campo magnetico esterno, $\vec{\sigma}$ vettore delle matrici di Pauli.
		\item \textbf{Evoluzione unitaria esplicita:}
		$U(t) = e^{-iHt/\hbar} = \cos(\omega_L t/2)I - i\sin(\omega_L t/2)\hat{n}\cdot\vec{\sigma}$
		con $\omega_L = \gamma|\vec{B}|$ (frequenza di Larmor).
		\item \textbf{Lato classico --- rotazione su $S^2$:} il valore di aspettazione
		$\langle\vec{\sigma}\rangle(t) = \mathrm{tr}(\rho(t)\vec{\sigma})$ ruota attorno a
		$\vec{B}$ con frequenza $\omega_L$; questo è esattamente il flusso hamiltoniano
		generato da $f_H(\theta,\phi) = \langle H\rangle$ su $S^2$ (Cap.\ 3).
		\item \textbf{Conclusione:} la precessione di Larmor è lo \emph{stesso} flusso
		hamiltoniano, realizzato in due modi: come rotazione su $S^2$ (lato classico,
		$C^\infty(S^2)$) e come evoluzione unitaria su $\mathbb{C}^2$ (lato quantistico,
		$M_2(\mathbb{C})$). Il teorema di Stone li unifica.
	\end{itemize}
	
	% ============================================================
	\newpage
	\section{Geometria Simplettica dello Spazio degli Stati}
	% ============================================================
	
	\nota{Questo capitolo offre una seconda lettura della risposta alla domanda della tesi,
		dal lato degli \emph{stati} invece che delle \emph{osservabili}. La risposta alla
		domanda è già completa al Cap.\ 7; questo capitolo la illumina geometricamente:
		lo spazio degli stati quantistico $\mathbb{P}\mathcal{H}$ è esso stesso una varietà
		simplettica (kähleriana), chiudendo il cerchio con il Cap.\ 1. Le strutture classiche
		non ``entrano'' nella meccanica quantistica dall'esterno --- sono già lì, nella
		geometria intrinseca dello spazio degli stati. La trattazione del qubit (§\ref{sec:qubit6})
		è completamente elementare: tutto si svolge in algebra lineare su $\mathbb{C}^2$,
		senza geometria differenziale complessa.}
	
	% ----------------------------------------------------------
	\subsection{Cornice Esplicita: dal Lato delle Osservabili a Quello degli Stati}
	% ----------------------------------------------------------
	
	\fonti{Ashtekar \& Schilling, \textit{Geometrical Formulation of Quantum Mechanics};
		Cirelli, Lanzavecchia \& Mania.}
	
	\begin{itemize}
		\item \textbf{Punto di vista complementare:} i Capp.\ 4--7 hanno mostrato che
		le strutture classiche sono struttura interna delle $C^*$-algebre delle osservabili.
		Questo capitolo guarda lo stesso risultato dallo spazio degli stati: anche
		$\mathbb{P}\mathcal{H}$ porta le strutture del Cap.\ 1.
		\item \textbf{Ridondanza di fase:} $|\psi\rangle$ e $e^{i\theta}|\psi\rangle$
		producono gli stessi valori di aspettazione; lo spazio fisico degli stati puri è
		il quoziente $\mathbb{P}\mathcal{H} = (\mathcal{H} \setminus \{0\})/\sim$.
	\end{itemize}
	
	% ----------------------------------------------------------
	\subsection{Decomposizione del Prodotto Scalare}
	% ----------------------------------------------------------
	
	\fonti{Ashtekar \& Schilling, \textit{Geometrical Formulation of Quantum Mechanics};
		Cirelli, Lanzavecchia \& Mania.}
	
	\begin{itemize}
		\item \textbf{Decomposizione fondamentale:} su $\mathcal{H}$ visto come spazio
		vettoriale reale, il prodotto scalare si scrive
		\[
		\langle\psi|\phi\rangle = g(\psi,\phi) + i\,\omega(\psi,\phi)
		\]
		dove $g(\psi,\phi) = \mathrm{Re}\langle\psi|\phi\rangle$ è una metrica riemanniana
		(simmetrica, definita positiva) e $\omega(\psi,\phi) = \mathrm{Im}\langle\psi|\phi\rangle$
		è una 2-forma antisimmetrica.
		\item \textbf{Non degenerazione di $\omega$:} $\omega(\psi,\phi) = 0$ per ogni
		$\phi$ implica $\psi = 0$; $\omega$ è non degenere, dunque simplettica.
		\item \textbf{Struttura complessa $J$:} la moltiplicazione $\psi \mapsto i\psi$
		è un endomorfismo reale di $\mathcal{H}$ con $J^2 = -\mathrm{Id}$; soddisfa
		$g(J\psi, J\phi) = g(\psi,\phi)$ e $\omega(\psi,\phi) = g(J\psi,\phi)$.
		\item \textbf{Struttura di Kähler su $\mathbb{P}\mathcal{H}$:} la terna
		$(g, J, \omega)$ discende al quoziente proiettivo e soddisfa le condizioni di
		compatibilità kähleriana; $\omega$ è chiusa (si verifica tramite il potenziale
		di Kähler $K = \log(1 + \sum|z_j|^2)$).
		\item \textbf{Nessuna geometria differenziale complessa necessaria:} tutto si
		riduce alla decomposizione in parte reale e parte immaginaria del prodotto scalare;
		la struttura kähleriana è intrinseca in $\mathcal{H}$, non imposta dall'esterno.
	\end{itemize}
	
	% ----------------------------------------------------------
	\subsection{Equazione di Schrödinger come Equazione Hamiltoniana}
	% ----------------------------------------------------------
	
	\fonti{Ashtekar \& Schilling, \textit{Geometrical Formulation of Quantum Mechanics}.}
	
	\begin{itemize}
		\item \textbf{Osservabile come funzione liscia su $\mathbb{P}\mathcal{H}$:}
		dato $\hat{A}$ autoaggiunto, $f_A([\psi]) = \langle\psi|\hat{A}|\psi\rangle /
		\langle\psi|\psi\rangle$ è liscia su $\mathbb{P}\mathcal{H}$; è l'analogo della
		funzione osservabile classica.
		\item \textbf{Equazione di Schrödinger in forma hamiltoniana:} le curve integrali
		del campo hamiltoniano $X_{f_H}$ definito da $\iota_{X_{f_H}}\omega = df_H$
		soddisfano $i\hbar\partial_t|\psi\rangle = \hat{H}|\psi\rangle$; la struttura
		simplettica su $\mathbb{P}\mathcal{H}$ è il dato che genera la dinamica,
		non un assioma aggiunto.
		\item \textbf{Identità esplicita:} le parentesi di Poisson su $\mathbb{P}\mathcal{H}$
		soddisfano $\{f_A, f_B\}_\omega = \frac{1}{i\hbar}\langle[\hat{A},\hat{B}]\rangle$;
		non un'analogia --- un'identità geometrica.
	\end{itemize}
	
	% ----------------------------------------------------------
	\subsection{Qubit VI: $\mathbb{C}^2$ in Coordinate Reali}
	\label{sec:qubit6}
	% ----------------------------------------------------------
	
	\fonti{Ashtekar \& Schilling, \textit{Geometrical Formulation of Quantum Mechanics}.}
	
	\begin{itemize}
		\item \textbf{Coordinate reali:} $\psi = (x_1 + iy_1, x_2 + iy_2)^T \in \mathbb{C}^2
		\cong \mathbb{R}^4$ con coordinate $(x_1, y_1, x_2, y_2)$.
		\item \textbf{Forma simplettica esplicita:}
		\[
		\omega = dx_1 \wedge dy_1 + dx_2 \wedge dy_2;
		\]
		in forma matriciale $\Omega = \mathrm{diag}(J_2, J_2)$ con $J_2 = \begin{pmatrix}
			0 & 1 \\ -1 & 0\end{pmatrix}$. Calcolo esplicito in coordinate.
		\item \textbf{Precessione di Larmor come flusso hamiltoniano:} l'equazione di
		Schrödinger del qubit con $H = -\gamma B\sigma_z/2$ si riscrive come
		$\dot{\mathbf{x}} = \Omega^{-1}\nabla f_H(\mathbf{x})$ in $\mathbb{R}^4$;
		le equazioni di Hamilton su $(\mathbb{R}^4, \omega)$ con hamiltoniana $f_H$.
		\item \textbf{Chiusura del cerchio:} la struttura simplettica con cui la tesi
		apre (Cap.\ 1) è la stessa struttura che governa lo spazio degli stati del
		sistema quantistico più semplice; le strutture classiche non sono entrate
		nella meccanica quantistica --- erano già qui.
	\end{itemize}
	
	% ============================================================
	\newpage
	\section{Noether, Taglio di Heisenberg e Risposta Sintetica}
	% ============================================================
	
	\nota{Il capitolo conclusivo svolge tre funzioni in sequenza: chiude il teorema di Noether
		(rimasto implicito fino qui); introduce il taglio di Heisenberg come questione filosofica
		ben posta nel linguaggio della tesi (non un confine ontologico, ma un limite asintotico
		nella norma C*); formula la risposta sintetica. Le tre funzioni sono sequenziali:
		Noether si applica al linguaggio del Cap.\ 3 e del Cap.\ 9; il taglio si contestualizza
		tramite il campo continuo di Rieffel del Cap.\ 6; la risposta sintetica riassume i tre
		passi dell'Introduzione.}
	
	% ----------------------------------------------------------
	\subsection{Teorema di Noether Quantistico}
	% ----------------------------------------------------------
	
	\fonti{Cirelli, Lanzavecchia \& Mania; Ashtekar \& Schilling.}
	
	\begin{itemize}
		\item \textbf{Enunciato:} se $[\hat{H}, \hat{A}] = 0$, allora
		$\frac{d}{dt}\langle\hat{A}\rangle_\psi = 0$ lungo ogni soluzione dell'equazione
		di Schrödinger.
		\item \textbf{Dimostrazione geometrica su $\mathbb{P}\mathcal{H}$:} $[\hat{H}, \hat{A}]
		= 0$ implica $\{f_H, f_A\}_\omega = 0$ (Cap.\ 9); per il teorema di Noether del
		Cap.\ 3 applicato a $(\mathbb{P}\mathcal{H}, \omega)$, $f_A$ è conservata lungo
		il flusso di $X_{f_H}$.
		\item \textbf{Stessa struttura del caso classico:} il generatore della simmetria
		commuta (in senso di Lie-bracket) con l'hamiltoniana; la conservazione segue
		per identità di Jacobi. Classico e quantistico usano la stessa dimostrazione ---
		perché usano la stessa struttura algebrica.
	\end{itemize}
	
	% ----------------------------------------------------------
	\subsection{Il Taglio di Heisenberg nel Linguaggio di Rieffel}
	% ----------------------------------------------------------
	
	\fonti{Landsman, \textit{Mathematical Topics Between Classical and Quantum Mechanics};
		Zurek, \textit{Decoherence and the Transition from Quantum to Classical}
		(Phys.\ Today, 1991).}
	
	\begin{itemize}
		\item \textbf{Il problema di Copenhagen:} il taglio di Heisenberg separa sistema
		quantistico e apparato di misura classico; Copenhagen lo postula senza localizzarlo;
		il confine è necessario ma arbitrario.
		\item \textbf{Il linguaggio del campo continuo:} se $\{A_\hbar\}$ è continua in
		$\hbar$, il confine classico-quantistico non è una discontinuità ontologica ---
		è la scelta di quale fibra $A_\hbar$ usare per descrivere il sistema.
		\item \textbf{Il taglio come numero:} nel modello del qubit (Cap.\ 6, §6.5),
		la decoerenza porta $\rho(t) \to \rho_{\mathrm{cl}}$ con
		$\|\rho(t) - \rho_{\mathrm{cl}}\|_1 = 2|c|e^{-\Gamma t}$; il taglio di
		Heisenberg è la condizione $\Gamma t \gg 1$ --- non un confine ontologico, è
		un numero.
		\item \textbf{Decoerenza come meccanismo fisico:} l'accoppiamento con un ambiente
		sopprime esponenzialmente le coerenze off-diagonali; matematicamente, sposta
		il sistema verso la fibra classica $A_0 = C_0(M)$ nel campo continuo di Rieffel.
	\end{itemize}
	
	% ----------------------------------------------------------
	\subsection{Risposta Sintetica alla Domanda della Tesi}
	% ----------------------------------------------------------
	
	\fonti{Landsman, \textit{Mathematical Topics Between Classical and Quantum Mechanics};
		Ashtekar \& Schilling; Strocchi, \textit{An Introduction to the Mathematical
			Structure of Quantum Mechanics}.}
	
	\begin{itemize}
		\item \textbf{Le strutture non riemergono --- non se ne sono mai andate.}
		Meccanica razionale e meccanica quantistica sono fibre dello stesso campo continuo
		di $C^*$-algebre (Rieffel): $A_0 = C_0(M)$ (fibra classica) e $A_\hbar$
		(fibra quantistica). Le strutture --- Lie-bracket, flusso hamiltoniano, Noether
		--- sono struttura dell'algebra, non della fisica specifica.
		\item \textbf{Il Lie-bracket:} è interno a ogni $C^*$-algebra; il commutatore
		e le parentesi di Poisson sono la stessa struttura in due rappresentazioni diverse
		(commutativa e non commutativa) della stessa algebra di Lie.
		\item \textbf{Il flusso hamiltoniano:} il teorema di Stone nel caso quantistico
		è l'analogo esatto del teorema di esistenza del flusso hamiltoniano nel caso
		classico; entrambi generano il flusso dal Lie-bracket con l'hamiltoniana.
		\item \textbf{Il teorema di Noether:} vale in entrambi i casi perché dipende
		solo dalla struttura simplettica e dall'identità di Jacobi --- strutture presenti
		in $(\mathbb{P}\mathcal{H}, \omega)$ come in $(M, \omega)$.
		\item \textbf{La domanda riformulata:} la domanda corretta non è ``perché le
		strutture classiche riemergono nel quantistico?'' ma ``perché ci aspettavamo
		che non ci fossero?''. La risposta è che la premessa della domanda era sbagliata:
		le due teorie non sono mai state separate da un confine strutturale ---
		solo da un valore del parametro $\hbar$.
	\end{itemize}
	
	% ============================================================
	\newpage
	\section*{Appendici}
	\addcontentsline{toc}{section}{Appendici}
	% ============================================================
	
	\subsection*{Appendice A --- Algebre di Lie}
	
	Raccoglie le definizioni e i risultati di teoria delle algebre di Lie usati
	nel testo: definizione assiomatica, omomorfismi, rappresentazioni, algebra
	inviluppante universale. Serve come riferimento autonomo per il lettore che
	non ha seguito un corso di algebra di Lie.
	
	\subsection*{Appendice B --- Teoria Spettrale}
	
	Richiami di teoria spettrale per operatori non limitati: operatori simmetrici e
	autoaggiunti, teorema spettrale (versione per operatori non limitati), calcolo
	funzionale misurabile. Necessaria per il Cap.\ 5 (teorema di Stone) e il Cap.\ 8
	(equazione di Schrödinger).
	
	\subsection*{Appendice C --- $C^*$-Algebre: Complementi}
	
	Completamenti tecnici del Cap.\ 4 che avrebbero appesantito il testo principale:
	dimostrazioni selezionate del teorema GNS, costruzione esplicita dell'algebra di
	Weyl per $n > 1$, enunciato del teorema di Fell (equivalenza tra rappresentazioni
	debolmente equivalenti). Necessaria per il lettore che vuole verificare le
	costruzioni del Cap.\ 6.
	
\end{document}